\documentclass[11pt]{article}

\usepackage{amsmath,amssymb}
\usepackage{xurl}
\usepackage[hidelinks]{hyperref}
\setlength{\emergencystretch}{2em}

% Shared commands for notes in docs/paper-gaps/.
%
% Two halves.  The link commands below are project identity: OWNER/REPO is the
% placeholder that scripts/bootstrap_project.py replaces with this project's
% GitHub slug and Pages site.  Everything after them is NOTATION, and notation
% belongs to the source paper: the definitions here are an example set, and a
% project replaces them with the macros of its own paper's style file.  The one
% rule that never changes: a command defined here must mean exactly what the
% source paper's macro means, or a note silently says something else.

\newcommand{\Lean}{\textsc{Lean}}
\newcommand{\leanid}[1]{\textsf{#1}}
\newcommand{\ghissue}[1]{\href{https://github.com/OWNER/REPO/issues/#1}{GitHub issue \##1}}
\newcommand{\ghpr}[1]{\href{https://github.com/OWNER/REPO/pull/#1}{PR \##1}}
% Local issue tracker (issues/NNNN-*.md in this repository); no remote URL.
\newcommand{\localissue}[1]{issue \##1 (\path{issues/})}
\newcommand{\doclink}[1]{\href{https://OWNER.github.io/REPO/docs/#1.html}{\path{#1}}}

% --- Notation: replace with the source paper's own macros --------------------
% \F, \N, \C, \R, \tr, \ot, \eps, \E, \bx, \bu, \bv, \by

\newcommand{\F}{\mathbb{F}}
\newcommand{\N}{\mathbb{N}}
\newcommand{\C}{\mathbb{C}}
\newcommand{\R}{\mathbb{R}}
\newcommand{\tr}{\mathrm{tr}}
\newcommand{\eps}{\epsilon}
\newcommand{\ot}{\otimes}
\newcommand{\E}{\mathop{\bf E\/}}

% bold random variables
\newcommand{\bu}{\boldsymbol{u}}
\newcommand{\bv}{\boldsymbol{v}}
\newcommand{\bx}{{\boldsymbol{x}}}
\newcommand{\by}{\boldsymbol{y}}

% T1 so literal < and > in \texttt placeholders typeset as themselves
\usepackage[T1]{fontenc}

% bra-ket
\usepackage{braket}

% --- Example structured spaces ---
\newcommand{\polyfunc}[3]{\mathcal{P}(#1, #2, #3)}
\newcommand{\polymeas}[3]{\mathrm{PolyMeas}(#1, #2, #3)}
\newcommand{\polysub}[3]{\mathrm{PolySub}(#1, #2, #3)}

% Approximation relations (paper uses \approx_\eps and \simeq_\zeta)
\newcommand{\approxeps}{\approx_{\varepsilon}}
\newcommand{\simgeq}{\simeq_{\zeta}}


\newcommand{\ind}{\operatorname{ind}}
\renewcommand{\ghissue}[1]{%
  \href{https://github.com/Dengnifer/MIPStarRE-A/issues/#1}{GitHub issue \##1}}

\title{The Probability of Anticommuting Pauli Tuples}
\author{}
\date{2026-09-04}

\begin{document}
\maketitle

\section{At a glance}

\begin{itemize}
  \item \textbf{Difficulty.} Low.  A product expansion computes the
  probability exactly and exposes two false claims in the printed proof.
  \item \textbf{Estimated weight.} The formal correction consists of one
  finite-field counting theorem and elementary real inequalities, roughly
  100--180 lines.
  \item \textbf{Mathlib/project split.} About 40\% finite counting and real
  inequalities, and 60\% project-local finite-field coordinates and sampling.
  \item \textbf{Key Mathlib inputs.} Uniform finite distributions, field-trace
  fibers, Bernoulli's inequality, and the fixed self-dual basis.  The formal
  proof is tracked in \ghissue{110}.
\end{itemize}

\section{Key theorem forms}

Let $q=2^k$, let $u_X,u_Z$ be uniform in $\F_q^m$, and let $r_X,r_Z$ be
uniform in $\F_q$.  Put
\[
  \gamma=\tr\big((r_X\ind_m(u_X))\cdot(r_Z\ind_m(u_Z))\big).
\]
Then
\[
  \Pr[\gamma\ne0]=\frac12(1-q^{-1})^{m+1},
  \qquad
  \Pr[\gamma=0]\geq\frac12.
\]
In particular, $m\leq q$ gives the uniform lower bound
$\Pr[\gamma\ne0]\geq2^{-4}$.  On the source parameter domain, where
$m,d\geq1$, the two lower bounds printed in the source also hold:
\[
  \Pr[\gamma\ne0]
  \geq \frac12(1-q^{-m})(1-q^{-1})(1-md/q)
  \geq \frac12(1-3md/q),
\]
and $\Pr[\gamma=0]\geq\frac12(1-3md/q)$.

\section{The source assertion and discrepancy}

Fact \path{fact:omega-anticomm-prob} of the Pauli analysis
\cite[lines~64--93]{Ji2020MIPStar} states the last two bounds without the
positivity assumptions on $m$ and $d$.  Its proof starts by
asserting that $\ind_m(u_X)=0$ if and only if $u_X=0$, and that $u_X=0$ forces
$\gamma=0$.%
\footnote{See
\path{references/qpbt-paper/14_analysis_of_the_pauli_basis_test.tex},
lines~64--93.}
Both assertions are false, and both are refuted inside the source parameter
domain.  The coordinates of $\ind_m(u)$ sum to one, so the indicator vector
never vanishes.  And for $(q,m)=(2,1)$ and $(u_X,u_Z,r_X,r_Z)=(0,0,1,1)$ the
phase is nonzero although $u_X=0$, so vanishing of $u_X$ does not force
$\gamma=0$; the degree $d$ does not occur in $\gamma$, so this holds for every
$d$, in particular for $d=1$.

The displayed bounds themselves are nevertheless correct throughout that
domain.  The source declares $\N$ to be the set of positive integers and takes
the parameter tuple $(q,m,d)$ in $\N$,%
\footnote{See \path{references/qpbt-paper/04_preliminaries.tex}, line~6, and
\path{references/qpbt-paper/08_classical_and_quantum_low_degree_tests.tex},
lines~903--904.  The blueprint records the same domain, and its admissibility
predicate repeats $d\geq1$ explicitly.}
so $m,d\geq1$ hold by hypothesis and the calculation of the next section
derives both printed bounds from them.  No source hypothesis is therefore
missing here: the printed proof is wrong, but its stated conclusion is not.
The bounds do fail outside the source domain --- at $(q,m,d)=(2,1,0)$ the exact
probability is $1/8$, below the printed $1/2$ --- which is why the formal
statements, whose parameters range over the Lean natural numbers and so include
zero, carry $m,d\geq1$ as an explicit hypothesis.  That hypothesis is a
faithful encoding of the source domain in the sense of
\path{docs/paper-gaps/proof-gap-protocol.tex}, not a missing source hypothesis
and not proof debt.

\section{The corrected calculation}

In characteristic two, expansion over the Boolean cube gives
\[
  \ind_m(u_X)\cdot\ind_m(u_Z)
  =\prod_{i=1}^m(1+u_{X,i}+u_{Z,i}).
\]
For fixed $u_X$, the factors are independent and uniform as $u_Z$ varies, so
they are all nonzero with probability $(1-q^{-1})^m$.  Independently,
$r_X\ne0$ with probability $1-q^{-1}$.  Conditioned on these events, the
argument of the trace is uniform as $r_Z$ varies, and a nonzero linear trace
to $\F_2$ takes the value one with probability $1/2$.  This proves the exact
formula.  Bernoulli's inequality gives the stated source bounds whenever
$m,d\geq1$, which both the source parameter domain and admissibility of the
Pauli-test parameters supply.  The role of $d\geq1$ is exactly that it forces
$md/q\geq m/q$, so that
$(1-q^{-1})^m\geq1-m/q\geq(1-q^{-m})(1-md/q)$, which is the first printed
bound; the second follows from it because $q^{-m}\leq q^{-1}\leq md/q$.

\section{Blueprint and formalization status}

The blueprint states the exact probability, restores $m,d\geq1$ on the
source bounds, and uses the unconditional constant lower bound when
conditioning the winning implications.  Its proof and correction remark cite
this note.  The corresponding formal declarations have the same corrected
statements, but their proofs remain open in the current tree.%
\footnote{The declarations are \leanid{anticommProb\_eq},
\leanid{commProb\_ge\_half}, \leanid{anticommProb\_ge\_of\_m\_le\_q}, and
\leanid{anticommProb\_ge\_of\_one\_le\_md} in
\doclink{MIPStarRE/QPBT/Observables/Anticommuting.lean}.}

\section{Conclusion}

The source proof is incorrect: both of its opening assertions are refuted
inside the source parameter domain.  Its displayed bounds are nevertheless true
throughout that domain, where $m,d\geq1$ by the source's convention that $\N$
is the set of positive integers, so no source hypothesis is missing.  The exact
probability is $\frac12(1-q^{-1})^{m+1}$; it yields universal conditioning
probabilities for all admissible parameters, and it recovers the printed lower
bounds once the positivity built into the source domain is written out --- as
the formal statements must do, since the Lean natural numbers include zero.

\bibliographystyle{alpha}
\bibliography{references}

\end{document}
